%% =====================================================================
%%  MOCKUP: Worked-Example Refined-Index Report   (m022)
%%
%%  All topology / matrices / series below are REAL (computed from
%%  manifold_index.core for m022).  Prose around the Dehn-filled series
%%  is marked [notional] where the point is layout, not numeric value.
%%
%%  LaTeX-notation conventions follow the in-app formatters:
%%    index_fmt.format_series_latex  (eta^{±nW_a}, q^{k/2})
%%    filling_fmt_v2.format_refined_index_notation
%%       → single-cusp:  I(e·α − (m/2)·β)  with no subscript on α,β
%%       → multi-cusp:   I(e₁α₁ − (m₁/2)β₁ + …)
%%       → Dehn filling: NC basis (γ, δ), slope transform γ-δ coords
%%
%%  Compile:  latexmk -pdf report_mockup_m022.tex
%% =====================================================================

\documentclass[11pt,a4paper]{article}

\usepackage[margin=1in]{geometry}
\usepackage{amsmath,amssymb,amsthm}
\usepackage{mathtools}
\usepackage{booktabs}
\usepackage{longtable}
\usepackage{array}
\usepackage{xcolor}
\usepackage{pifont}
\usepackage{hyperref}
\hypersetup{colorlinks=true, linkcolor=black, urlcolor=blue}

%% ── Shorthand macros (mirroring the in-app formatters) ─────────────
\newcommand{\I}{\mathcal{I}}                          % bare index symbol
\renewcommand{\Im}{\mathcal{I}_{m022}}                % refined index, manifold-subscripted
                                                      %   (overrides amssymb's \Im = Im(z))
\newcommand{\IDelta}{\mathcal{I}_{\Delta}}            % tetrahedron index
%  Filled index: subscript is (manifold_name)_{user slope in α,β}
%  Usage: \Ifill{m022}{5\alpha+2\beta}
\newcommand{\Ifill}[2]{\mathcal{I}_{(#1)_{#2}}}
\newcommand{\Zhalf}{\tfrac{1}{2}\mathbb{Z}}
\newcommand{\pass}{{\color{green!55!black}\ding{51}}}
\newcommand{\fail}{{\color{red!70!black}\ding{55}}}
\newcommand{\tet}[1]{\Delta_{#1}}

\title{Refined 3D Index --- Worked Example \\[2pt] \large \texttt{m022}}
\author{Generated by \texttt{manifold-index}}
\date{\today}

\begin{document}
\maketitle
\tableofcontents
\bigskip

%% =====================================================================
\section{Setup}
%% =====================================================================

We walk through the full refined-index computation for the census
manifold \texttt{m022}.  Every intermediate quantity below is the value
actually produced by \texttt{manifold\_index.core}; the layout is
pedagogical so a reader can reproduce any step by hand.

\begin{center}
\begin{tabular}{l l}
\toprule
Manifold                        & \texttt{m022} \\
Volume                          & \(2.98912\ldots\) \\
Tetrahedra \(n\)                & \(4\) \\
Cusps \(r\)                     & \(1\) \\
Internal edges \(n-r\)          & \(3\) \;(\(3\) easy, \(0\) hard) \\
Series order                    & \(q^{20/2}\) \;(\(N_{\max}=10\)) \\
Display sectors (per cusp)      & \((\alpha,\beta) \in
   \bigl\{(0,0),\;(0,-\tfrac12),\;(\tfrac12,0),\;(0,-1),\;(1,0)\bigr\}\) \\
\bottomrule
\end{tabular}
\end{center}

\paragraph{Notation.}
Every displayed index carries the manifold name as a subscript, as in
the in-app series table.  An unfilled sector reads
\(\Im(e\alpha - \tfrac{m}{2}\beta)\) with \(\alpha\) the meridian,
\(\beta\) the longitude, and \((m,e)\) the internal charge variables
related by \((\alpha,\beta) = (e,\,-m/2)\).  Hard edges (if present)
carry fugacities \(\eta^{nW_a}\) (\(0\)-indexed in \(a\)); residual
cusp \(\eta\)'s carry \(\eta^{nV_i}\).  The tetrahedron factor is
\(\IDelta(m_a, e_a)\).  A Dehn-filled index along user slope
\(P\alpha+Q\beta\) is written
\(\Ifill{m022}{P\alpha+Q\beta}\) --- the manifold name moves
inside a parenthesised subscript to make room for the slope, matching
\texttt{build\_fill\_row\_cells}.

\paragraph{A note on \texttt{m022}.}
The edge classification (Step 2) returns \(0\) hard edges: every
internal edge equation can already be solved for some \(\eta\).  That
means \texttt{m022} carries \emph{no} \(\eta\)-fugacities, and the
``refined'' index coincides with the ordinary \(3\)-D index.  The
Weyl condition (Step 6) reduces to the \(m\mapsto-m\) symmetry of the
pure \(q\)-series.  For a manifold with non-empty hard-edge set, the
sectors would additionally carry \(\eta^{\pm nW_a}\) structure.

%% =====================================================================
\section{Step 1 --- Triangulation data}
%% =====================================================================

SnapPy reports the gluing equations as a \((n+2r)\times 3n = 6\times 12\)
integer matrix.  Each row encodes exponents \((f_i,g_i,h_i)_{i=1}^{n}\)
of \((z_i, z_i', z_i'')\) in a relation
\(\prod_i z_i^{f_i}\, z_i'^{g_i}\, z_i''^{h_i} = 1\).

\paragraph{Edge rows \((n=4)\).}  Columns group in triples
\((f_i, g_i, h_i)\) for \(i=1,\dots,4\).
\[
G_{\text{edge}} = \scriptsize
\left(\begin{array}{rrr|rrr|rrr|rrr}
 2 & 0 & 0 & 1 & 0 & 0 & 2 & 0 & 0 & 0 & 0 & 1 \\
 0 & 2 & 1 & 0 & 0 & 2 & 0 & 2 & 1 & 0 & 2 & 0 \\
 0 & 0 & 1 & 0 & 2 & 0 & 0 & 0 & 0 & 0 & 0 & 1 \\
 0 & 0 & 0 & 1 & 0 & 0 & 0 & 0 & 1 & 2 & 0 & 0 \\
\end{array}\right)
\]

\paragraph{Cusp rows \((2r=2)\): meridian \(\alpha\), longitude \(\beta\).}
\[
G_{\text{cusp}} = \scriptsize
\left(\begin{array}{rrr|rrr|rrr|rrr}
 2 & -1 & 0 & 1 & 0 & -1 & 0 & -1 & -1 & 0 & -1 & 0 \\
-1 &  1 & 1 & -1 & -1 & 1 & 0 & 2 & 0 & 0 & 2 & 0 \\
\end{array}\right)
\]

\emph{How to read:} the four edge rows state that each ideal edge glues
up to total holonomy \(1\); the two cusp rows give the products that
express meridian and longitude around the single cusp.

%% =====================================================================
\section{Step 2 --- Edge classification}
%% =====================================================================

An edge row is \emph{easy} if, after column-wise inspection, at most one
of \((z_i, z_i', z_i'')\) is non-zero per tetrahedron; otherwise
\emph{hard}.  For \texttt{m022}:

\begin{center}
\begin{tabular}{c c l}
\toprule
Label       & Independent? & Row (original index)  \\
\midrule
E\(_0\)     & \pass        & edge 0 \\
E\(_1\)     & \pass        & edge 3 \\
E\(_2\)     & \pass        & edge 2 \\
\bottomrule
\end{tabular}
\end{center}

No hard edges are needed, so the hard-padding list is empty; the
\(n-r=3\) internal coordinates \(e_{\text{int},0},
e_{\text{int},1}, e_{\text{int},2}\) correspond directly to the three
independent easy edges \(\{E_0, E_1, E_2\}\).

%% =====================================================================
\section{Step 3 --- Neumann--Zagier data}
%% =====================================================================

Symplectic reduction produces
\(g_{\text{NZ}} \in \mathrm{Sp}(2n, \mathbb{Q})\) of shape \(8\times 8\).
Columns are ordered \((Z_1, Z_2, Z_3, Z_4,\; Z_1'', Z_2'', Z_3'', Z_4'')\);
rows are grouped
\(\bigl[\,\text{meridians}_r \,\vert\, \text{hard}_{d_h} \,\vert\,
   \text{easy}_{d_e} \,\vert\, \text{longitudes/2}_r \,\vert\,
   \text{int.\ momenta}_{n-r}\,\bigr]\),
here \([\,1 \,\vert\, 0 \,\vert\, 3 \,\vert\, 1 \,\vert\, 3\,]\).

\[
g_{\text{NZ}} =
\left(\begin{array}{rrrr|rrrr}
 3 &  1 &  1 &  1 &  1 & -1 &  0 &  1 \\
 2 &  1 &  2 &  0 &  0 &  0 &  0 &  1 \\
 0 &  1 &  0 &  2 &  0 &  0 &  1 &  0 \\
 0 & -2 &  0 &  0 &  1 & -2 &  0 &  1 \\ \hline
-1 &  0 & -1 & -1 &  0 &  1 & -1 & -1 \\
-2 & -1 & -1 & -2 & -1 &  1 &  0 & -1 \\
-3 & -1 & -2 & -1 & -1 &  1 &  0 & -1 \\
-2 &  0 & -1 & -1 & -1 &  1 &  0 & -1 \\
\end{array}\right)
\]

\paragraph{Affine shifts.}
\[
\nu_x = (-3,\; -2,\; -2,\; 0), \qquad
\nu_p = ( 2,\; 0,\; 0,\; 0).
\]

\emph{Interpretation:} \(\nu_p\) shifts the overall \((-q^{1/2})\) exponent;
\(\nu_x\) shifts the \(\eta\)-monomial (trivially, here, since there are
no \(\eta\)).  Both are intrinsic to the triangulation.

%% =====================================================================
\section{Step 4 --- Applying the formula to one sector}
%% =====================================================================

Before listing the full sector grid, walk through the single sector
\((m_{\text{ext}}, e_{\text{ext}}) = (1, 0)\) --- the display label
\(\Im\!\left(0\alpha - \tfrac{1}{2}\beta\right)\) --- in full.

\paragraph{(a) Build \(\kappa\).}
\(\kappa = (m_{\text{ext}},\; \mathbf{0}^{n-r},\;
           e_{\text{ext}},\; \mathbf{e}_{\text{int}})
 = (1, 0, 0, 0,\; 0, e_0, e_1, e_2)\),
with \(\mathbf{e}_{\text{int}} = (e_0, e_1, e_2) \in (\Zhalf)^3\) to be
summed over.

\paragraph{(b) Local charges via \(g_{\text{NZ}}^{-1}\kappa\).}
\[
g_{\text{NZ}}^{-1} =
\left(\begin{array}{rrrr|rrrr}
 0 & -1 & -1 & -1 & -1 &  0 &  0 & -1 \\
 1 &  1 &  1 &  1 &  1 &  0 &  0 &  2 \\
-1 &  0 &  0 &  0 &  0 &  0 & -1 &  0 \\
-1 & -1 & -1 & -1 & -1 & -1 &  0 & -1 \\ \hline
 1 &  2 &  3 &  2 &  3 &  2 &  0 &  0 \\
 0 &  1 &  1 &  0 &  1 &  1 &  1 & -2 \\
 1 &  1 &  2 &  1 &  1 &  2 &  0 &  0 \\
 1 &  2 &  1 &  1 &  1 &  0 &  2 &  0 \\
\end{array}\right)
\]
Reading off rows \(1{:}n\) and \(n{+}1{:}2n\) and specialising to
\(m_{\text{ext}}=1,\; e_{\text{ext}}=0\):
\begin{align*}
\tet{1}:\quad & m_1 = -\,e_2, & e_1 &= 1 + 2e_0, \\
\tet{2}:\quad & m_2 = 1 + 2e_2, & e_2 &= e_0 + e_1 - 2e_2, \\
\tet{3}:\quad & m_3 = -1 - e_1, & e_3 &= 1 + 2e_0, \\
\tet{4}:\quad & m_4 = -1 - e_0 - e_2, & e_4 &= 1 + 2e_1.
\end{align*}

\paragraph{(c) Integrality constraints.}
All \(m_a, e_a\) must be integers.  Since \(2e_k\) is always integer
for \(e_k\in\Zhalf\), the \(e_a\) rows are automatically integer.
The \(m_a\) rows force
\(e_2 \in \mathbb{Z},\; e_1 \in \mathbb{Z},\;
  e_0 + e_2 \in \mathbb{Z}\),
hence \(\mathbf{e}_{\text{int}} \in \mathbb{Z}^3\) (not a strict
half-step).  The internal sum therefore runs over ordinary integer
lattice points rather than the full \((\Zhalf)^3\).

\paragraph{(d) Prefactor and result.}
Each admissible \(\mathbf{e}_{\text{int}}\) contributes
\[
(-q^{1/2})^{\mathbf{m}\cdot\nu_p \,-\, \mathbf{e}\cdot\nu_x}
\prod_{a=1}^{4} \IDelta(m_a, e_a).
\]
Summing over \(\mathbf{e}_{\text{int}} \in \mathbb{Z}^3\) and truncating
at \(q^{10}\):
\[
\Im\!\left(0\alpha - \tfrac{1}{2}\beta\right) \;=\;
  -\,q \;+\; 3\,q^{3} \;+\; 5\,q^{4} \;+\; q^{5}
  \;-\; 16\,q^{6} \;-\; 37\,q^{7} \;-\; 57\,q^{8}
  \;-\; 53\,q^{9} \;+\; 9\,q^{10}
  \;+\; O(q^{11}).
\]

%% =====================================================================
\section{Step 5 --- All computed sectors}
%% =====================================================================

The app displays a canonical grid of \(5^r\) sectors per cusp.  For
\texttt{m022} (\(r=1\)) the grid has \(5\) entries, of which \(4\) are
non-zero.  Truncation: \(q^{20/2}\), \(N_{\max}=10\).  In each row the
argument on the left is
\(e\alpha - \tfrac{m}{2}\beta\) as in the in-app index table.

\begin{align*}
\Im\!\left(0\alpha + 0\beta\right) &=
  1 - q - q^{2} + 4q^{3} + 6q^{4} + 4q^{5}
  - 16q^{6} - 41q^{7} - 68q^{8} - 73q^{9}
  - 9q^{10} + O(q^{11}), \\[2pt]
\Im\!\left(0\alpha - \tfrac{1}{2}\beta\right) &=
  -q + 3q^{3} + 5q^{4} + q^{5} - 16q^{6}
  - 37q^{7} - 57q^{8} - 53q^{9} + 9q^{10} + O(q^{11}), \\[2pt]
\Im\!\left(\tfrac{1}{2}\alpha + 0\beta\right) &= 0, \\[2pt]
\Im\!\left(0\alpha - \beta\right) &=
  2q^{3} + q^{4} - 2q^{5} - 14q^{6} - 24q^{7}
  - 28q^{8} - 14q^{9} + 47q^{10} + O(q^{11}), \\[2pt]
\Im\!\left(\alpha + 0\beta\right) &=
  -q + 3q^{3} + 5q^{4} + q^{5} - 14q^{6}
  - 35q^{7} - 54q^{8} - 49q^{9} + 9q^{10} + O(q^{11}).
\end{align*}

\emph{Observation:} Because the sector argument is linear in the signed
charges \((e,-m/2)\), the Weyl condition
\(\Im(e\alpha-\tfrac{m}{2}\beta) = \Im(-e\alpha+\tfrac{m}{2}\beta)\)
is the statement that flipping the sign of every component gives the
same \(q\)-series.  With no hard edges, there is no accompanying
\(\eta^{\pm nW_a}\) shift to cancel.

%% =====================================================================
\section{Step 6 --- Weyl-symmetry check}
\label{sec:weyl}
%% =====================================================================

For a general manifold the Weyl-shifted index
\[
f(\eta_j; m, e) \;=\;
   \eta_j^{\sum_I a_{j,I}e_I + b_{j,I}m_I}\;
   \I_{\!X}\!\bigl(e\alpha - \tfrac{m}{2}\beta\bigr)
\]
is required to satisfy \(f(m,e) = f(-m,-e)\) for each hard edge \(j\),
together with a pair of integrality conditions on \((a_j, b_j)\) and an
adjoint-character projection fixed to \(-1\) per filled cusp.

For \texttt{m022} the hard-edge set is empty, so the Weyl vectors
\(\mathbf{a}, \mathbf{b}\) are \(0\)-dimensional:
\begin{center}
\begin{tabular}{c c c c c}
\toprule
\(j\) & \(a_j\) & \(b_j\) & \(a_j\in\mathbb{Z}\)? &
                    \(2b_j\in\mathbb{Z}\)? \\
\midrule
--- & --- & --- & \pass\,(vacuous) & \pass\,(vacuous) \\
\bottomrule
\end{tabular}
\end{center}

\paragraph{Sector-pair symmetry.}
\(4/4\) non-zero sectors in the \(5\)-entry display grid satisfy
\(\Im(e\alpha-\tfrac{m}{2}\beta) = \Im(-e\alpha+\tfrac{m}{2}\beta)\)
\pass.

\paragraph{Adjoint projection.}
With no \(\eta\) integration, the adjoint-character step reduces to
\(\sum_e c_e\,\chi_{\text{adj}}(e) = -1\), where \(c_e\) is the
\(q^{1}\) coefficient of \(\Im\!\left(e\alpha + 0\beta\right)\).
For \texttt{m022}:
\begin{center}
\begin{tabular}{c r}
\toprule
\(e\) (\(=\alpha\)-coeff) & \(c_e\) \\
\midrule
\(-\tfrac12\) & \(0\) \\
\(0\)         & \(-1\) \\
\(+\tfrac12\) & \(0\) \\
\bottomrule
\end{tabular}
\end{center}
Projected value: \(-1\) \;(expected \(-1\)) \;\pass.

%% =====================================================================
\section{Step 7 --- Dehn filling, worked example}
%% =====================================================================

Fill the single cusp along the user slope \(P_{\text{user}}\alpha +
Q_{\text{user}}\beta\) with \((P_{\text{user}}, Q_{\text{user}}) = (5, 2)\).

\subsection*{(a) Non-closable cycle search}
Search range \(P\in[-10,10],\; Q\in[1,10]\) \;[notional range].  One
NC cycle returned:
\[
\gamma \;=\; P_{\text{nc}}\,\alpha + Q_{\text{nc}}\,\beta, \qquad
(P_{\text{nc}}, Q_{\text{nc}}) = (1, 1). \quad \text{[notional]}
\]

\subsection*{(b) Complementary basis vector \(\delta\)}
Solve \(R\,Q_{\text{nc}} - P_{\text{nc}}\,S = 1\) by extended Euclid:
\((R, S) = (0, -1)\), giving \(\delta = 0\cdot\alpha + (-1)\cdot\beta\).
Verify
\(\det\!\begin{pmatrix}P_{\text{nc}}&Q_{\text{nc}}\\R&S\end{pmatrix}
  = (1)(-1) - (1)(0) = -1 \ne 0\) \pass.

\subsection*{(c) User slope in the \((\gamma,\delta)\) basis}
Expand \(P_{\text{user}}\alpha+Q_{\text{user}}\beta = p\gamma + q\delta\):
\[
\begin{pmatrix}p\\q\end{pmatrix}
=
\begin{pmatrix} S & -R \\ -Q_{\text{nc}} & P_{\text{nc}} \end{pmatrix}
\begin{pmatrix} P_{\text{user}} \\ Q_{\text{user}} \end{pmatrix}
=
\begin{pmatrix} -5 \\ -3 \end{pmatrix}.
\]
The filled index is displayed with the user slope (not the transformed
one) in the subscript, as
\(\Ifill{\texttt{m022}}{5\alpha + 2\beta}\); the transformed slope
\(p/q = (-5)/(-3)\) enters the HJ step next.

\subsection*{(d) Hirzebruch--Jung continued fraction}
\(p/q = (-5)/(-3)\).  The HJ algorithm returns \([2, 2, 2]\), length
\(\ell = 3\). \;[notional; \texttt{hj\_continued\_fraction} would be run
on real \((p, q)\)]

\subsection*{(e) Basis change on \(g_{\text{NZ}}\)}
\texttt{apply\_cusp\_basis\_change} replaces the \((\alpha, \beta/2)\)
rows of \(g_{\text{NZ}}\) with
\[
\text{new-position} = P_{\text{nc}}\,\alpha + 2Q_{\text{nc}}\,(\beta/2),
\qquad
\text{new-momentum} = a\,\alpha + b\,(\beta/2),
\]
where \((a,b)\) satisfy the Bézout identity
\(P_{\text{nc}}\,b - 2Q_{\text{nc}}\,a = 1\).  The result:
\[
g_{\text{NZ}}' \;=\;
\left(\text{[\,8$\times$8 integer matrix, printed in full in the
real report\,]}\right),\quad
\nu_x' = (\text{updated}),\quad \nu_p' = (\text{updated}).
\]

\subsection*{(f) Physical Weyl vectors and edge compatibility}
With \(0\) hard edges, \(\mathbf{a}_{\text{phys}}\) and
\(\mathbf{b}_{\text{phys}}\) are empty \(0\)-vectors; no edge is
collapsed; all hard edges (vacuously) integer-compatible. \pass

\subsection*{(g) Kernel product \& contraction}
HJ kernel assembly for \([2,2,2]\) produces a truncated polynomial with
\(\sim 17\) non-zero terms at the chosen \(q\)-order.  Contracting
against the pre-computed \(\Im\) and collecting residual cusp-\(\eta\)
fugacities \(\eta^{nV_0}\):
\[
\Ifill{\texttt{m022}}{5\alpha + 2\beta} \;=\;
  1 + \bigl(-2 - \eta^{2V_0} - \eta^{-2V_0}\bigr)q
    + \bigl(\eta^{4V_0} + \eta^{-4V_0} - 4\bigr)q^{2} + \cdots + O(q^{11}).
\quad \text{[notional coefficients]}
\]
(For a manifold with hard edges, terms of the form
\(\eta^{\pm nW_a}\eta^{nV_0}\) would also appear; cf.\ the
\((5_2)_{\alpha+2\beta}\) series in the in-app display.)

%% =====================================================================
\section*{Appendix --- Tetrahedron index definition}
%% =====================================================================
\addcontentsline{toc}{section}{Appendix --- Tetrahedron index definition}

\[
I_t(m, e) \;=\; \sum_{k=\max(0,-e)}^{\infty}
  \frac{(-1)^k \, q^{k(k+1)/2 - (2k+e)\,m/2}}
       {\prod_{j=1}^{k}(1-q^j)\;\prod_{j=1}^{k+e}(1-q^j)},
\qquad
\IDelta(m, e) \;=\;
\begin{cases}
(-q^{1/2})^m \, I_t(-m-e,\; m), & m+e \ge 0, \\[2pt]
I_t(m, e),                      & m+e < 0.
\end{cases}
\]

\end{document}
